\chapter{力学}
本章では運動を決める力、保存量、物質の変形を結び付ける。物理学で既習の運動方程式を再導入するより、試験で落としやすい条件の区別を明確にする。力は加速度を決め、エネルギーは始状態と終状態を結び、応力は物体内部に力がどう分配されるかを表す。速度と加速度、静止摩擦と動摩擦、浮力と合力、半径と高度を区別すれば、同じ式を異なる年度にも適用できる。

\section{力と加速度：斜面・摩擦・見掛けの重さ}
\label{theory:mechanics:newton}
慣性系では $m\ddot{\boldsymbol r}=\sum\boldsymbol F$ である。合力がゼロなら速度は一定であり、必ずしも静止していない。鉛直投射の最高点では速度が一瞬ゼロになるが、重力による加速度はゼロにならない。この区別が\hyperref[q:r04-reg-applied-04]{令和4年度問4}の核心である。

斜面の上向きを正とし、引く力 $F$ が斜面に平行なら $N=mg\cos\theta$ である。上向きに滑り始める限界では摩擦が下向きに働くため
\[
 F=mg\sin\theta+\mu_s mg\cos\theta.
\]
力を加えずに下向きに滑り始める角度 $\theta_c$ が分かれば、$mg\sin\theta_c=\mu_s mg\cos\theta_c$ より $\mu_s=\tan\theta_c$ と決まる。静止摩擦は常に $\mu_sN$ なのではなく、その値以下でつり合いに必要な大きさを取る。滑り続ける物体では動摩擦係数 $\mu_k$ を使う。平成23年度問1は後者、平成28・令和3・令和7年度問1は前者である。

体重計は質量を直接測らず、足から受ける力を標準の重力加速度で換算している。上向き加速度 $a$ のエレベーターなら $N=m(g+a)$、一定速度なら $a=0$ なので読みは変わらない。自由落下で読みがゼロになっても質量は消失しない。質量は力と加速度の関係や運動量からも測定できる。

\section{仕事とエネルギー：途中の経路を消去する}
\label{theory:mechanics:energy}
運動方程式と変位の内積を取ると $d(mv^2/2)=\boldsymbol F\cdot d\boldsymbol r$ となる。重力と理想ばねだけが仕事をする場合、
\[
 E=\frac12 mv^2+mgh+\frac12 kx^2
\]
が保存される。抵抗・摩擦による仕事がある場合には、この三項だけの保存を仮定しない。

高さ $h$ から初速の大きさ $v_0$ で投射すると、着地時の速さは $v^2=v_0^2+2gh$ である。投射方向は飛行時間や着地点を変えるが、同じ始終高度なら速さを変えない。令和元年度再試験問2はこの不変量を問う。平成24年度問1は $v_0=0$、令和4年度問1は水平ばねのエネルギー $kx^2/2$ を運動エネルギーへ変換する場合である。数値を代入する前に質量が消えるかを確かめれば、余計な計算を避けられる。

\section{弾性：材料定数と形状による剛性を分ける}
\label{theory:mechanics:elasticity}
応力 $\sigma=F/A$ は単位面積当たりの内力、ひずみ $\epsilon=\Delta L/L$ は相対変形である。小変形・線形弾性の一軸引張では $\sigma=E\epsilon$ が成り立つ。これがフックの法則の材料表示であり、ヤング率 $E$ の単位は Pa である。一様棒のばね定数は $k=EA/L$ なので、形状で変わる剛性 $k$ と材料の $E$ を取り違えない。

横ひずみ $\epsilon_\perp$ と軸ひずみ $\epsilon_\parallel$ からポアソン比を $\nu=-\epsilon_\perp/\epsilon_\parallel$ と定義する。通常の正の $\nu$ を持つ材料は引張方向に伸び、横方向に縮む。一次の体積ひずみは三方向のひずみの和なので
\[
 \frac{\Delta V}{V}\simeq\epsilon_x+\epsilon_y+\epsilon_z=(1-2\nu)\epsilon_\parallel.
\]
したがって令和5年度問1で与えられた $0.1\leq\nu\leq0.45$ の材料は引張で体積が増加する。$\nu=1/2$ は非圧縮の極限である。すべての材料がこの数値範囲に入ると一般化しない。

せん断では形状の角度変化 $\gamma$ に対し $\tau=G\gamma$ であり、$G$ がせん断弾性率である。等方線形弾性では $E=2G(1+\nu)$ となる。名称の空欄を埋める平成23年度問4・平成25年度問6・平成30年度問2も、軸変形と角度変形を区別する問題である。これらの構成関係はMITの材料力学講義にも整理されている\cite{mit_torsion}。

\section{丸棒のねじり：最大せん断応力は表面に生じる}
\label{theory:mechanics:torsion}
一様な丸棒の線形弾性ねじりでは、中心から距離 $r$ のせん断ひずみは $r\,d\phi/dz$、応力は $\tau=Gr\,d\phi/dz$ である。断面全体のトルクを積分すると
\[
 T=\int_A\tau r\,dA=G\frac{d\phi}{dz}J,
 \qquad J=\int_A r^2\,dA=\frac{\pi d^4}{32}.
\]
よって $\tau=Tr/J$、表面 $r=d/2$ では $\tau_{\max}=16T/(\pi d^3)$ となる。$J$ は断面二次極モーメントで、単位は $\mathrm{m^4}$ である。令和元年度問2では mm を m に直してから三乗する。直径を半径として代入すると8倍の誤差になる。この導出と円形断面・線形弾性という適用条件はMIT講義のねじり節を参照されたい\cite{mit_torsion}。非円形断面や塑性化した棒へ同じ分布を無条件に適用しない。

\section{静水圧と浮力：置き換えた流体の重さ}
\label{theory:mechanics:hydrostatics}
静止流体の鉛直上向き座標を $z$ とすると、薄い流体要素の力のつり合いから $dp/dz=-\rho_f g$ となる。物体表面の圧力を積分した上向き合力は、同じ体積を占める流体の重さと等しい。密度一定なら
\[
 F_B=\rho_f gV_{\mathrm{displaced}}.
\]
完全に水没した物体では排除体積が物体全体の体積であり、浮力に物体密度は入らない。物体密度 $\rho_b$ が入るのは重力との合力 $(\rho_f-\rho_b)Vg$ や、浮いているときの水没割合 $V_{\mathrm{displaced}}/V=\rho_b/\rho_f$ である。平成26年度問3と令和元年度問1は同じ式を異なる番号の選択肢で問うため、番号ではなく式を記憶する。

\section{流れのエネルギーとピトー管}
\label{theory:mechanics:bernoulli}
定常・非粘性・非圧縮の流れで、同じ流線に沿ってオイラー方程式を積分すると
\[
 p+\frac12\rho v^2+\rho gz=\mathrm{const.}
\]
となる。各項は単位体積当たりのエネルギーであり、圧力の単位を持つ。水平流では静圧 $p$ と動圧 $q=\rho v^2/2$ の和が全圧 $p_0$ である。流れを正面の孔で停止させるピトー管は全圧を、流れに沿う側孔は静圧を測り、差から $v=\sqrt{2(p_0-p)/\rho}$ を求める。この測定配置はMITの実験講義資料で確認できる\cite{mit_pitot}。

平成24年度問9・令和2年度問2では「全圧と静圧を測定し、その差が動圧」という関係を問う。粘性抵抗を扱うストークスの法則と混同しない。大きな圧縮性、圧力損失、異なる流線間では、この簡単な式の条件を再点検する必要がある。

\section{円軌道：第一宇宙速度と静止衛星}
\label{theory:mechanics:orbit}
万有引力が円運動の向心力になるので $GMm/r^2=mv^2/r$、すなわち $v=\sqrt{GM/r}$ である。地表近くで $g=GM/R^2$ とおけば第一宇宙速度は $\sqrt{gR}$ になる。衛星質量が消える点が平成26年度問1・平成28年度問2の識別点である。

角速度 $2\pi/T$ を使うと軌道半径は
\[
 r=\left(\frac{GMT^2}{4\pi^2}\right)^{1/3}.
\]
静止衛星は地球の自転と同じ角速度で赤道上を順行する。概算では $T\simeq24$ 時間を秒に換算する。地心距離 $r$ は約 $4.2\times10^7$ m、地表からの高度は $r-R$ であり別の量である。平成30年度問1は m、令和6年度問1は km で選択肢を示す。厳密には恒星日を使うが、提示された桁の選択には概算との差は影響しない。

\section{衝突と反発：保存則の適用区間を選ぶ}
\label{theory:mechanics:impact}
短い衝突時間では外力の力積を無視できる条件の下で、衝突する二物体の運動量が保存される。反発係数は離れる相対速度と近づく相対速度の比であり、一直線上なら $v_B-v_A=e(u_A-u_B)$ である。$e=1$ で運動エネルギーも保存されるが、$e<1$ では減少する。

質量 $m$、速度 $u$ の物体が静止する質量 $M$ の物体へ完全弾性衝突すると
\[
 v_A=\frac{m-M}{m+M}u,\qquad v_B=\frac{2m}{m+M}u.
\]
令和3年度問2では求める物体がBであることを確認する。地面で跳ね返る球なら、落下区間では $v^2=2gh$、衝突で速さが $e$ 倍、上昇区間で高さが $e^2$ 倍となる。2回目の反発高さは $e^4h_0$ であり、令和2年度問1の「2回」を1回の式へ代入しない。

\section{剛体のつり合い：支点を選んで未知力を減らす}
\label{theory:mechanics:equilibrium}
剛体の静止条件は合力と合モーメントがともにゼロである。床に立てた一様棒の長さを $L$、床となす角を $\theta$ とする。壁が滑らかなら壁の反力 $N_w$ は水平、床では鉛直反力 $N_f$ と摩擦力 $F$ が働く。床との接点を支点に選ぶと $N_f,F$ のモーメントはゼロになり、
\[
 N_wL\sin\theta=Mg\frac L2\cos\theta,
 \qquad F=N_w=\frac{Mg}{2\tan\theta},\qquad N_f=Mg
\]
と決まる。令和6年度問2では $L$ が消えることを次元解析でも確認できる。ただしこの静止状態の実現には $F\leq\mu_sN_f$ が必要である。必要な摩擦力を求める式と、摩擦限界で倒れる条件は別である。

\section{振動：復元力の傾きが周期を決める}
\label{theory:mechanics:oscillation}
平衡点からの変位 $x$ に対して復元力が $-kx$ なら $m\ddot x+kx=0$、$\omega=\sqrt{k/m}$、$T=2\pi\sqrt{m/k}$ である。鉛直ばねは静的伸び $\delta=mg/k$ を用いると $T=2\pi\sqrt{\delta/g}$ となる。重力は平衡点を移すが、線形ばねの復元力の傾きを変えない。平成27年度問1・令和7年度問2はこの形で整理できる。

単振り子は $mL^2\ddot\theta=-mgL\sin\theta$ であり、小振幅で $\sin\theta\simeq\theta$ と近似すれば $T=2\pi\sqrt{L/g}$ となる。質量は消える。平成29年度問1・令和5年度問2の短い数値計算はこの近似を使う。大振幅では周期が振幅に依存するため、任意振幅への一般化はしない。基礎科目の\hyperref[q:h24-reg-basic-3-2]{平成24年度III-2}および\hyperref[q:h26-reg-basic-3-5]{平成26年度III-5}とも同じ線形振動の構造を共有する。

\section*{原典確認・独立解法を保存した問題}
以下の解説は査読待ちの初稿である。章の共通理論へ戻り、年度ごとの条件と各選択肢を検討する。正答番号の一致のみを最終検証とはしない。
\subsection{平成23年度 問1}
\label{q:h23-reg-applied-01}

\noindent\textbf{原典：} sources/H23/17\_applied\_science.pdf、PDF物理1頁。

\textbf{公式正答：} 3（single）。

\par\textit{原典本文の校正と独自解説は未作成。分類のみの登録項目。}

\subsection{平成23年度 問4}
\label{q:h23-reg-applied-04}

\noindent\textbf{原典：} sources/H23/17\_applied\_science.pdf、PDF物理2頁。

\textbf{公式正答：} 1（single）。

\par\textit{原典本文の校正と独自解説は未作成。分類のみの登録項目。}

\subsection{平成24年度 問1}
\label{q:h24-reg-applied-01}

\noindent\textbf{原典：} sources/H24/17\_applied\_science.pdf、PDF物理1頁。

\textbf{公式正答：} 3（single）。

\par\textit{原典本文の校正と独自解説は未作成。分類のみの登録項目。}

\subsection{平成24年度 問9}
\label{q:h24-reg-applied-09}

\noindent\textbf{原典：} sources/H24/17\_applied\_science.pdf、PDF物理3頁。

\textbf{公式正答：} 1（single）。

\par\textit{原典本文の校正と独自解説は未作成。分類のみの登録項目。}

\subsection{平成25年度 問1}
\label{q:h25-reg-applied-01}

\noindent\textbf{原典：} sources/H25/17\_applied\_science.pdf、PDF物理1頁。

\textbf{公式正答：} 2（single）。

\par\textit{原典本文の校正と独自解説は未作成。分類のみの登録項目。}

\subsection{平成25年度 問6}
\label{q:h25-reg-applied-06}

\noindent\textbf{原典：} sources/H25/17\_applied\_science.pdf、PDF物理3頁。

\textbf{公式正答：} 1（single）。

\par\textit{原典本文の校正と独自解説は未作成。分類のみの登録項目。}

\subsection{平成26年度 問1}
\label{q:h26-reg-applied-01}

\noindent\textbf{原典：} sources/H26/17\_applied\_science.pdf、PDF物理1頁。

\textbf{公式正答：} 1（single）。

\par\textit{原典本文の校正と独自解説は未作成。分類のみの登録項目。}

\subsection{平成26年度 問3}
\label{q:h26-reg-applied-03}

\noindent\textbf{原典：} sources/H26/17\_applied\_science.pdf、PDF物理1頁。

\textbf{公式正答：} 2（single）。

\par\textit{原典本文の校正と独自解説は未作成。分類のみの登録項目。}

\subsection{平成28年度 問1}
\label{q:h28-reg-applied-01}

\noindent\textbf{原典：} sources/H28/17\_applied\_science.pdf、PDF物理1頁。

\textbf{公式正答：} 5（single）。

\par\textit{原典本文の校正と独自解説は未作成。分類のみの登録項目。}

\subsection{平成28年度 問2}
\label{q:h28-reg-applied-02}

\noindent\textbf{原典：} sources/H28/17\_applied\_science.pdf、PDF物理1頁。

\textbf{公式正答：} 1（single）。

\par\textit{原典本文の校正と独自解説は未作成。分類のみの登録項目。}

\subsection{平成30年度 問1}
\label{q:h30-reg-applied-01}

\noindent\textbf{原典：} sources/H30/17\_applied\_science.pdf、PDF物理1頁。

\textbf{公式正答：} 4（single）。

\par\textit{原典本文の校正と独自解説は未作成。分類のみの登録項目。}

\subsection{平成30年度 問2}
\label{q:h30-reg-applied-02}

\noindent\textbf{原典：} sources/H30/17\_applied\_science.pdf、PDF物理1頁。

\textbf{公式正答：} 2（single）。

\par\textit{原典本文の校正と独自解説は未作成。分類のみの登録項目。}

\subsection{令和元年度 問1}
\label{q:r01-reg-applied-01}

\noindent\textbf{原典：} sources/R01/17\_applied\_science.pdf、PDF物理1頁。

\textbf{公式正答：} 5（single）。

\par\textit{原典本文の校正と独自解説は未作成。分類のみの登録項目。}

\subsection{令和元年度 問2}
\label{q:r01-reg-applied-02}

\noindent\textbf{原典：} sources/R01/17\_applied\_science.pdf、PDF物理1頁。

\textbf{公式正答：} 2（single）。

\par\textit{原典本文の校正と独自解説は未作成。分類のみの登録項目。}

\subsection{令和元年度再試験 問2}
\label{q:r01-ret-applied-02}

\noindent\textbf{原典：} sources/R01\_retest/17\_applied\_science.pdf、PDF物理1頁。

\textbf{公式正答：} 4（single）。

\par\textit{原典本文の校正と独自解説は未作成。分類のみの登録項目。}

\subsection{令和2年度 問1}
\label{q:r02-reg-applied-01}

\noindent\textbf{原典：} sources/R02/17\_applied\_science.pdf、PDF物理1頁。

\textbf{公式正答：} 2（single）。

\par\textit{原典本文の校正と独自解説は未作成。分類のみの登録項目。}

\subsection{令和2年度 問2}
\label{q:r02-reg-applied-02}

\noindent\textbf{原典：} sources/R02/17\_applied\_science.pdf、PDF物理1頁。

\textbf{公式正答：} 1（single）。

\par\textit{原典本文の校正と独自解説は未作成。分類のみの登録項目。}

\subsection{令和3年度 問1}
\label{q:r03-reg-applied-01}

\noindent\textbf{原典：} sources/R03/17\_applied\_science.pdf、PDF物理1頁。

\textbf{公式正答：} 4（single）。

\par\textit{原典本文の校正と独自解説は未作成。分類のみの登録項目。}

\subsection{令和3年度 問2}
\label{q:r03-reg-applied-02}

\noindent\textbf{原典：} sources/R03/17\_applied\_science.pdf、PDF物理1頁。

\textbf{公式正答：} 5（single）。

\par\textit{原典本文の校正と独自解説は未作成。分類のみの登録項目。}

\subsection{令和4年度 問1}
\label{q:r04-reg-applied-01}

\noindent\textbf{原典：} sources/R04/17\_applied\_science.pdf、PDF物理1頁。

\textbf{公式正答：} 3（single）。

\par\textit{原典本文の校正と独自解説は未作成。分類のみの登録項目。}

\subsection{令和4年度 問2}
\label{q:r04-reg-applied-02}

\noindent\textbf{原典：} sources/R04/17\_applied\_science.pdf、PDF物理1頁。

\textbf{公式正答：} 5（single）。

\par\textit{原典本文の校正と独自解説は未作成。分類のみの登録項目。}

\subsection{令和5年度 問1}
\label{q:r05-reg-applied-01}

\noindent\textbf{原典：} sources/R05/17\_applied\_science.pdf、PDF物理1頁。

\textbf{公式正答：} 5（single）。

\par\textit{原典本文の校正と独自解説は未作成。分類のみの登録項目。}

\subsection{令和5年度 問2}
\label{q:r05-reg-applied-02}

\noindent\textbf{原典：} sources/R05/17\_applied\_science.pdf、PDF物理1頁。

\textbf{公式正答：} 1（single）。

\par\textit{原典本文の校正と独自解説は未作成。分類のみの登録項目。}

\subsection{令和6年度 問1}
\label{q:r06-reg-applied-01}

\noindent\textbf{原典：} sources/R06/17\_applied\_science.pdf、PDF物理1頁。

\textbf{公式正答：} 3（single）。

\par\textit{原典本文の校正と独自解説は未作成。分類のみの登録項目。}

\subsection{令和6年度 問2}
\label{q:r06-reg-applied-02}

\noindent\textbf{原典：} sources/R06/17\_applied\_science.pdf、PDF物理1頁。

\textbf{公式正答：} 5（single）。

\par\textit{原典本文の校正と独自解説は未作成。分類のみの登録項目。}

\subsection{令和7年度 問1}
\label{q:r07-reg-applied-01}

\noindent\textbf{原典：} sources/R07/17\_applied\_science.pdf、PDF物理1頁。

\textbf{公式正答：} 3（single）。

\par\textit{原典本文の校正と独自解説は未作成。分類のみの登録項目。}

